% from ega4-2-fr.tex
\oldpage[IV-2]{68}

En particulier les composantes irréductibles (resp. connexes) de $X$ sont
géométriquement irréductibles (resp. géométriquement connexes).
\end{proposition}

\begin{proof}
Par définition, $k'$ est une extension radicielle de $k$~; le morphisme
$\operatorname{Spec}(k')\to\operatorname{Spec}(k)$ est entier, surjectif et
radiciel, donc la première assertion résulte de (2.4.5, (i))~; la seconde en
résulte, vu (4.5.1).
\end{proof}

\subsection{Préschémas algébriques géométriquement réduits}
\label{subsection:IV.2.4.6-fr}

\begin{proposition}[4.6.1]
\label{IV.2.4.6.1-fr}
Soient $k$ un corps, $X$ un $k$-préschéma, $\Omega$ une extension parfaite de
$k$. Les conditions suivantes sont équivalentes~:
\begin{enumerate}
\item[a)] Pour tout $k$-préschéma réduit $S$, $X\times_k S$ est réduit.
\item[b)] Pour toute extension $K$ de $k$, $X\otimes_k K$ est réduit.
\item[c)] $X\otimes_k\Omega$ est réduit.
\item[d)] Pour toute extension radicielle finie $k'$ de $k$, $X\otimes_k k'$ est
réduit.
\item[e)] $X$ est réduite, et pour toute composante irréductible $X_\alpha$ de
$X$, de point générique $x_\alpha$, $k(x_\alpha)$ est une extension séparable de
$k$.
\end{enumerate}
\end{proposition}

\begin{proof}
Il est trivial que a) $\Rightarrow$ b) $\Rightarrow$ c)~; comme $k'$ peut être
considérée comme sous-extension de $\Omega$, on a vu (4.4.1) que c) implique
d). Montrons que d) implique e). En prenant $k'=k$, on voit d'abord que $X$
est réduit~; pour prouver que $k(x_\alpha)$ est séparable sur $k$, il suffit de
montrer que $k(x_\alpha)\otimes_k k'$ est un anneau réduit pour toute extension
finie radicielle $k'$ de $k$ (4.3.5). On peut se borner au cas où
$X=\operatorname{Spec}(A)$ est affine, $A$ étant une $k$-algèbre réduite~; alors
par hypothèse l'anneau $A\otimes_k k'$ est réduit (I, 5.1.4) et
$k(x_\alpha)\otimes_k k'$ est un anneau de fractions de cet anneau, donc il est
réduit $(0_{\mathrm I},1.2.8)$. Enfin, pour prouver que e) entraîne a), on peut
se borner au cas où $X=\operatorname{Spec}(A)$, $S=\operatorname{Spec}(B)$ sont
affines, $A$ et $B$ étant des $k$-algèbres. Les $k(x_\alpha)$ sont les corps des
fractions des anneaux quotients $A/\mathfrak p_\alpha$, où
$\mathfrak p_\alpha=\mathfrak j_{x_\alpha}$ sont les idéaux premiers minimaux de
$A$~; comme par hypothèse $\bigcap_\alpha\mathfrak p_\alpha=0$, $A$ est contenu
dans le produit $\prod_\alpha A/\mathfrak p_\alpha$, donc dans
$\prod_\alpha k(x_\alpha)$~; on peut de même considérer $B$ comme une
sous-algèbre d'un produit $\prod_\beta L_\beta$ d'extensions de $k$. Il en
résulte que $A\otimes_k B$ s'identifie à une sous-algèbre de
\[
\left(\prod_\alpha k(x_\alpha)\right)\otimes_k
\left(\prod_\beta L_\beta\right),
\]
et ce produit tensoriel s'identifie lui-même à une sous-algèbre de
\[
\prod_{\alpha,\beta}\bigl(k(x_\alpha)\otimes_k L_\beta\bigr)
\]
(Bourbaki, \emph{Alg.}, chap.~II, 3\textsuperscript{e} éd., \S~7,
n\textsuperscript{o}~7, prop.~15). Or par hypothèse les
$k(x_\alpha)\otimes_k L_\beta$ sont réduits (4.3.5), donc il en est de même de
$A\otimes_k B$.
\end{proof}

\begin{definition}[4.6.2]
\label{IV.2.4.6.2-fr}
Si les conditions équivalentes a) à e) de la proposition (4.6.1) sont
remplies, on dit que $X$ est \emph{séparable} (ou \emph{géométriquement réduit},
ou \emph{universellement réduit}) sur $k$. On dit qu'un $k$-préschéma $X$ est
\emph{géométriquement intègre} sur $k$ si, pour toute extension $K$ de $k$,
$X\otimes_k K$ est intègre~; il revient au même (en vertu de (4.6.1)) de dire
que $X$ est séparable et géométriquement irréductible sur $k$.
\end{definition}

On dira qu'une $k$-algèbre $A$ (commutative) est \emph{séparable} si
$\operatorname{Spec}(A)$ est séparable sur $k$~; cela signifie donc que pour
toute extension $K$ de $k$, l'anneau $A_{(K)}=A\otimes_k K$ est réduit. On
notera que cette définition coïncide avec celle de Bourbaki, \emph{Alg.},
chap.~VIII,

